% ---------------------------------------------------------------------
%  PHỤ LỤC
% ---------------------------------------------------------------------
\chapter{Phụ lục}
\markboth{Phụ lục}{}

\section*{A. Hướng dẫn cài đặt và chạy chương trình}
\phantomsection
\addcontentsline{toc}{section}{A. Hướng dẫn cài đặt và chạy chương trình}

Để cài đặt và triển khai hệ thống FoodAI trên máy tính cá nhân hoặc máy chủ thử nghiệm, người dùng thực hiện tuần tự theo các bước dưới đây:

\begin{lstlisting}[language=bash, caption={Các bước cấu hình môi trường và khởi chạy ứng dụng FoodAI.}]
# 1. Clone ma nguon du an tu kho luu tru Git
git clone https://github.com/phamtanhao/food-recommendation-ai.git
cd food-recommendation-ai

# 2. Cau hinh bien moi truong cho Backend
cd backend
cp .env.example .env
# Chinh sua cac tham so trong file .env:
# GEMINI_API_KEY=AIzaSy... (Khoa Google Gemini API)
# PEXELS_API_KEY=...      (Khoa Pexels API)
# JWT_SECRET=super_secret_jwt_key_2026
# PORT=5000

# 3. Cai dat thu vien va khoi tao co so du lieu MySQL / SQLite
npm install
npm run db:init      # Tu dong tao cac bang va nap tai khoan mac dinh
npm run db:seed      # Nap kho tri thuc 220 mon an Cookpad tu JSON

# 4. Khoi chay may chu Backend (Lang nghe tai cong 5000)
npm run dev

# 5. Mo cua so Terminal moi de cai dat va chay Frontend
cd ../frontend
npm install
npm run dev          # Khoi chay ung dung Vite tai http://localhost:5173

# 6. Tai khoan dang nhap thu nghiem (Demo Accounts):
# - Tai khoan Admin: admin@foodsystem.com    / Mat khau: Admin@123456
# - Tai khoan User:  user@foodsystem.com     / Mat khau: User@123456
\end{lstlisting}

\section*{B. Bảng tự kiểm trước khi nộp đồ án}
\phantomsection
\addcontentsline{toc}{section}{B. Bảng tự kiểm trước khi nộp đồ án}

Bảng đối chiếu mười tiêu chuẩn chất lượng theo Rubric đánh giá học phần AIP202:

\begin{tukiem}
  \item Sản phẩm chạy ổn định trên môi trường máy tính khác, không phụ thuộc cục bộ.
  \item Đã chạy thử nghiệm đầy đủ các luồng nghiệp vụ trên đúng môi trường trình chiếu bảo vệ.
  \item Mã nguồn được quản lý chuyên nghiệp bằng kho lưu trữ \code{git}, phân chia nhánh và commit rõ ràng.
  \item Có đầy đủ tệp hướng dẫn \code{README.md}, \code{package.json} và tài liệu mô tả API.
  \item Mọi hình vẽ, sơ đồ kiến trúc, bảng số liệu và giải thuật đều có chú thích, nhãn tham chiếu và được trích dẫn trong thân bài.
  \item Mọi mục tiêu đề ra ở Chương~\ref{ch:modau} đều được đối chiếu và giải trình kết quả cụ thể ở Chương~\ref{ch:ketqua}.
  \item Mọi số liệu thực nghiệm trong báo cáo được đo đạc từ hệ thống chạy thật, bảo đảm tính trung thực học thuật.
  \item Mọi tài liệu tham khảo trong \code{refs.bib} đều được trích dẫn ít nhất một lần theo chuẩn IEEE.
  \item Đã khai báo minh bạch phần ứng dụng mô hình Trí tuệ nhân tạo (Google Gemini) và cơ chế kiểm soát thuật toán.
  \item Đã kiểm tra toàn bộ bản in PDF, loại bỏ hoàn toàn các hộp gợi ý và khối văn mẫu.
\end{tukiem}

\section*{C. Nhận xét của giảng viên hướng dẫn}
\phantomsection
\addcontentsline{toc}{section}{C. Nhận xét của giảng viên hướng dẫn}

\khungnhanxet[7.5cm]

\vspace{0.8em}
\noindent\begin{minipage}[t]{0.5\textwidth}
\color{daunavy}
{\color{daumaroon}\bfseries Điểm số}

\vspace{0.7em}
\dongchamnhan{Bằng số:}

\vspace{0.9em}
\dongchamnhan{Bằng chữ:}
\end{minipage}\hfill
\okyten{\dauGVHD}{Giảng viên hướng dẫn}
